\scanpage{100}
\mspage{93}
% archive sheet number: 99
Soit, pour tout $i\in I$, $\widetilde C_i$ l'ens ordonné
\[
\widetilde C_i=\{\text{sous-catégories }C'_i\text{ de }C_i=\underline M^{\pi}/X_i\mid
\begin{aligned}
&\operatorname{card}\operatorname{Fl}(C'_i)\leq\pi,\\
&C'_i\text{ filtrante}
\end{aligned}\}.
\]
J'admets pour l'instant le

\uline{Lemme} Pour \deleted{\unsure{tout}} tout, $X\in\deleted{\unsure{E|M^{\pi}}}\inserted{M^{\deleted{\unsure{\pi|\ast}}}}$, la catégorie $\deleted{\unsure{\underline M^{\pi}/X|C_i}}\inserted{\underline M^{\pi}/X}$ est filtrante.
\inserted{(\unsure{Bien|Par} l'hypothèse (b), \unsure{en ce que|on a que|en que} \unsure{tt élém de|tt pt de|tt partie de|tt objet de} $M$ est de la forme \unsure{dite dans a)|dit dans a)|et de dans}) \unsure{comme} $X_i\in M^{\deleted{\unsure{\pi|\ast}}}$, \unsure{par|puis} $C_i$ est filtrant,}
% margin note written sideways
\margin{Utiliser l'hypothèse \deleted{\unsure{précisée|forte}} que les éléments de $\mathrm{Filt}^{\pi}$ sont \unsure{$\pi$-accessibles|$\pi$-accessible}.}

Il en résulte, \deleted{\unsure{donc}} la réunion des $C'_i$ \inserted{$\in\widetilde C_i$} est $C_i$; et l'ensemble ordonné $\widetilde C_i$ est filtrant, \inserted{et} toute partie \inserted{$(C'_{i,\alpha})_{\alpha\in A}$ filtrante et} de cardinal $\leq\pi$ de $\unsure{\widetilde C'_i|\widetilde C_i}$ admet une borne supérieure $C'_{iA}$, savoir la réunion des $C'_{i\alpha}$. $\widetilde C_i$ est donc un ens ordonné $\pi$-adapté. \unsure{On}

On peut considérer les ens. ordonnés $\widetilde C_i$, pour $i\in I$ variable, comme les fibres d'une application croissante d'ens ordonnés
\[
p:\widetilde C\longrightarrow C,
\]
en \deleted{disant} \unsure{écrivant}
\[
\left\{\begin{aligned}
&(C'_i\subset C_i)\leq(C'_j\subset C_j)\\
&\text{ssi }i\leq j,\text{ et le foncteur canonique}\\
&\text{de transition }C_i\longrightarrow C_j\text{ applique }C'_i\\
&\text{dans }C'_j.
\end{aligned}\right.
\]
On voit donc que $\widetilde C$ est lui aussi filtrant et grand devant $\pi$.
